\section{Introduction}
\label{sec:introduction}

Real-time industrial systems combine discrete control logic with
quantitative timing requirements. In railway control, communication
protocols, and aerospace scheduling
~\cite{basile2021unisig,lonn1997tdma,david2010herschel},
correctness depends not only on which actions occur, but also on whether
they occur within permitted time windows. The scope of this paper is
timed safety properties rather than unbounded liveness properties,
because deadlines, bounded waiting, alarm durations, and error-state
reachability describe finite executions whose violations are directly
observable during design and testing
~\cite{kolbl2019clock,vogel2022patterns}. Timed automata (TAs) model
such behavior with clocks, guards, resets, and invariants
~\cite{alur1994theory,bengtsson2004timed}; practical systems are often
modeled as networks of timed automata (NTA).

Clock behavior is central to timed safety verification. Clocks and their
constraints encode deadlines, timeouts, durations, and delays. An
incorrect numerical bound in a transition guard or location invariant can
make otherwise valid discrete behavior occur too early, too late, or for
an invalid duration
~\cite{kolbl2019clock,jiang2010heart,lonn1997tdma}. This paper studies
clock-bound repair for NTA models, modifying only numerical bounds in
guards and invariants while leaving properties, locations, transitions,
synchronizations, clock references, resets, and data updates unchanged.
Model checkers such as UPPAAL, Kronos, and opaal can return timed
diagnostic traces (TDTs), namely timed counterexamples reaching
violating states
~\cite{larsen1997uppaal,daws1996kronos,dalsgaard2011opaal}. A TDT
exposes a violating execution, but it does not identify a globally
acceptable clock-bound edit.

Repair is more difficult when an NTA must satisfy multiple properties
simultaneously. A bound change that blocks one TDT may leave another
violation feasible, violate a previously satisfied property, or
compromise the original functional behavior. In this paper, functional
behavior means the untimed discrete behavior characterized by functional
equivalence. Thus, local trace elimination is necessary but insufficient:
a candidate edit should be accepted only when it satisfies the full
property set and preserves the original functional behavior. A practical
repair method should therefore generate small and inspectable
candidate edits from local TDT evidence, but accept a candidate only
after full-property regression verification and an admissibility check
based on untimed functional equivalence~\cite{kolbl2019clock}.

Existing repair techniques provide an important foundation. TARTAR
~\cite{kolbl2019clock} introduced TDT-based clock-bound repair by
encoding a TDT into linear real arithmetic constraints, solving a
MaxSMT problem for small bound changes, and checking admissibility
through functional equivalence. Related work studies parametric timed
automata, clock-guard repair, robustness analysis, and SMT-based repair
~\cite{alur1993parametric,andre2019repairing,bendik2021robustness,jose2011bugassist}.
However, these workflows are primarily trace-local and single-property
in style. When repairs are generated property by property, the coupling
among properties, counterexamples, and repairable bounds is not used to
guide candidate generation, and a local edit may be rejected only after
full validation.

This paper presents MPTA-Repair, an automated framework for simultaneous
multi-property clock-bound repair of NTA models via iterative solving.
The key idea is to treat each violated property and its TDT as a
property-trace constraint, maintain a shared constraint pool, and use
relations among properties, counterexamples, and repairable bounds to
guide MaxSMT-based candidate generation. Each candidate is validated by
full-property regression verification and an admissibility check based on
untimed functional equivalence. If validation exposes new violations,
their TDTs are added to the pool and the solver is invoked again.

We evaluate MPTA-Repair on benchmarks derived from UPPAAL
\cite{behrmann2004tutorial} and the XtaBenchmarkSuite
\cite{farkas2018benchmarks}, together with an industrial dataset
constructed in this work. Faulty NTA instances are produced by controlled
clock-bound mutations. Compared with an adapted iterative TARTAR-style
baseline, MPTA-Repair improves the success rate from $72\%$ to $93\%$
on the benchmark dataset and from $20\%$ to $76\%$ on the industrial
dataset.

This paper makes three contributions. First, it formulates simultaneous
multi-property clock-bound repair for NTA models under full-property
satisfaction and preservation of untimed functional behavior. Second, it
introduces an iterative MaxSMT workflow based on a shared property-trace
constraint pool and relation-guided candidate generation. Third, it
evaluates the framework on benchmark and industrial faulty models,
showing that joint multi-property repair is more effective than an
iterative trace-local baseline in this setting.

The remainder of the paper is organized as follows.
Section~\ref{sec:background} introduces the notation and
accepted-repair condition. Section~\ref{sec:approach} presents
MPTA-Repair. Section~\ref{sec:evaluation} reports the experimental
evaluation. Section~\ref{sec:conclusion} concludes the paper.